\begin{abstract}
国际油价由供需基本面、金融条件和地缘风险共同驱动。为避免混淆预测偏差、事件异常与战争净因果
贡献，本文基于2010年1月至2026年6月多源数据，建立“预测与事件识别--结构冲击传导--政策
反事实--尾部风险优化”框架；前问的冲击与响应核依次进入后问，实现四问贯通。

\textbf{针对问题一，}在78个共同目标月上以扩展窗口比较六类预测器。不变预测基线的1、3、6个月
RMSE分别为0.1408、0.2660和0.3082，均为最低；以2026年2月70.887美元/桶为原点，三期点预测
均为70.887美元/桶。状态匹配事件研究得到$CAR[0,+1]=$\textbf{15.61\%}（经验$p=0.0244$），
航运受阻阶段条件波动率为战前的2.214倍；递归SVAR将3月36.45\%的实际收益中的32.23\%
识别为供给收缩贡献，总结构贡献处于历史98.4\%分位。上述结果表明事件附近存在异常，
但不等同于战争净因果份额。

\textbf{针对问题二，}将186个月的三类SVAR冲击输入局部投影，并以2000次移动块bootstrap构造
95\%联合置信带。一次标准差石油特定风险冲击使人民币原油成本当月上升
\textbf{8.832\%}，联合区间为$[6.888\%,10.776\%]$；PPI、CPI峰值分别为0.762\%、0.247\%，
工业增加值第12个月谷值为$-0.427\%$。后三者联合区间均跨零，故模型支持“成本
先行、物价滞后、工业活动可能承压”的顺序，但\textbf{尚无稳健的总体增长损失证据}。

\textbf{针对问题三，}以燃油分布滞后、面板局部投影和动态反事实评价政策。六国6个月零售燃油
传导率中位数为0.227；中国调价指数代理为0.373，因价格层不同不纳入正式排名。部分识别的临界
缩放系数中位数为0.590，95\%区间为$[0.436,0.834]$。若关闭2026年临时调控，6月PPI、CPI将
分别高2.596\%（$[1.453\%,3.633\%]$）和0.414\%（$[0.199\%,1.088\%]$），工业增加值差额
为$-0.469\%$且区间跨零；物价缓冲同时形成\textbf{1425元/吨延期缺口}。因此现有证据
支持条件情景下的局部物价缓冲，但不足以证明我国总体韧性全面优于对照国。

\textbf{针对问题四，}FHS--GJR--GARCH的20000条路径给出6个月Brent中位数71.76美元/桶、
90\%区间$[37.15,117.90]$美元/桶，期末和路径内曾超过历史95\%阈值的概率分别为5.56\%和9.65\%。
提出的\textbf{SAPR--CVaR三状态自适应调价模型}从1771组规则中筛得119组1\%容差Pareto解，
最终传导率为$(0.30,0.10,0.10)$。在2022年至2026年的49个隔离检验情景中，相对完全传导，
平均宏观损失、95\%尾部损失和调价波动分别\textbf{降低47.70\%、38.79\%和30.28\%}，代价是
平均延期缺口负担增加2.22\%；2000次联合块bootstrap中，相对三种预注册基线保持非支配的
概率为100\%。该“最优”仅限本文注册的规则族、六个月情景库和四目标准则。

\keywords{国际油价\quad{}结构VAR\quad{}局部投影\quad{}政策反事实\quad{}SAPR--CVaR}
\end{abstract}
